% !TEX root = ../tfg_etsiinf_NombreApellidos.tex
\chapter{Conclusiones y líneas futuras} \label{chp:conclusiones}

Este capítulo valora el grado de cumplimiento de los objetivos planteados en el
Capítulo~\ref{chp:intro} y resume las conclusiones técnicas y metodológicas del
trabajo. La discusión se apoya en la evidencia cuantitativa presentada en el
Capítulo~\ref{chp:resultados}.

\section{Conclusiones} \label{sct:conclusiones:conclusiones}

El objetivo general del trabajo, desarrollar un entorno de aprendizaje por
refuerzo sobre el simulador de multirrotores de Aerostack2 y utilizarlo para
entrenar y evaluar un controlador de velocidad, se ha alcanzado de forma
satisfactoria. La política final resuelve el 100\,\% de los 100 episodios del
protocolo de evaluación determinista sobre la tarea completa (objetivos
aleatorios a hasta 5 m con orientación final), con una distancia final media de
0,32 m.

El primer objetivo específico, la familiarización con Aerostack2 y su simulador,
se alcanzó de forma natural durante las primeras fases y quedó reflejado en la
identificación de los componentes de control, estimación y comportamiento
relevantes para el problema. El segundo, la construcción de la interfaz
Gymnasium sobre el simulador, se cumplió con un matiz importante que constituye
una de las lecciones del trabajo: la interfaz mínima fue rápida de construir,
pero convertirla en un entorno \emph{válido} para aprendizaje por refuerzo exigió
un esfuerzo muy superior al previsto, concentrado en el reinicio certificado de
episodios y en la verificación de accionabilidad.

El tercer objetivo, la validación experimental del ciclo de interacción, se
alcanzó mediante una estrategia de tres niveles (pruebas deterministas,
integración con el simulador y ejecuciones PPO acotadas) que acabó articulada en
una veintena de suites de validación automatizadas. La consecución del cuarto
objetivo, el entorno vectorizado, resultó la más difícil: el primer
entrenamiento con dos drones colapsó por una interacción no documentada entre el
paso vectorizado en \emph{lockstep} y la persistencia de referencias de velocidad
del controlador, y solo un diagnóstico sistemático con pruebas discriminatorias
permitió identificar la causa raíz (8,05 m de deriva en 20 s sin comandos) y
corregirla con un mecanismo de vigilancia de comandos inactivos. Tras la
corrección, el entrenamiento vectorizado duplica el muestreo efectivo y mejoró
la política final del 78\,\% al 100\,\% de éxito determinista.

El quinto objetivo, la definición de un problema de control progresivo, resultó
ser la clave del aprendizaje. Con la recompensa mínima y la tarea completa, la
exploración no encontraba señal de éxito suficiente (tasas planas del 11\,\% aun
con la infraestructura ya libre de fallos); el curriculum de distancia en tres
etapas con reutilización del modelo anterior elevó la tasa de éxito del 27\,\% al
65\,\% y finalmente al 91--99\,\% sobre la tarea original. El sexto objetivo, el
entrenamiento, la evaluación y la comparación con un controlador PID, se cumplió
mediante un protocolo de episodios sembrados idénticos para todos los
controladores. La comparación es honesta: el PID gana la tarea nominal en tiempo
y precisión de orientación, la política alcanza la paridad en tasa de éxito, y
la diferencia queda explicada y cuantificada. La extensión jerárquica planteada
como séptimo objetivo opcional no se abordó y se recoge como línea futura.

Más allá de los objetivos, el trabajo deja dos conclusiones metodológicas. La
primera es que en aprendizaje por refuerzo sobre simuladores robóticos reales la
validez física del episodio es la condición que gobierna todo lo demás: la mayor
parte del esfuerzo de ingeniería se invirtió en garantizar que cada transición
entregada al algoritmo representara de verdad la tarea, y los mecanismos
resultantes (reinicio certificado con verificación de accionabilidad, guards de
acción predictivos, supervisión con reanudación y reciclado preventivo) son
aportaciones reutilizables independientes de la política entrenada. La segunda
es el valor de los resultados negativos tratados como evidencia: los fallos de
Exp010, Exp011 y Exp017 produjeron los tres hallazgos más importantes del
trabajo (la degradación acumulativa del simulador con los reinicios de servicio,
la insuficiencia de los guards reactivos frente a la inercia, y la persistencia
indefinida de las referencias de velocidad), todos ellos diagnosticados con
causa raíz y cuantificación.

El código desarrollado queda publicado en los dos repositorios indicados en el
Capítulo~\ref{chp:implementacion}, con configuraciones, pruebas automatizadas y
artefactos de evaluación reproducibles mediante semillas.

\section{Líneas futuras} \label{sct:conclusiones:futuro}

La línea futura principal, ya planteada por el tutor durante el trabajo y ahora
motivada cuantitativamente por tres hallazgos independientes, es sustituir ROS 2
en el bucle de entrenamiento por \emph{bindings} directos sobre la dinámica del
simulador, con llamadas explícitas de la forma estado--acción--paso temporal. Un
simulador pausable eliminaría de raíz la degradación por teleports (el estado se
asignaría directamente), la deriva por referencias persistentes (no habría
tiempo real entre pasos) y el coste del \emph{lockstep} vectorizado (los
reinicios serían instantáneos), y permitiría escalar la vectorización a decenas
de agentes aprovechando la inferencia por lotes. ROS 2 quedaría reservado para el
despliegue de la política entrenada a la frecuencia correspondiente al paso
temporal usado en entrenamiento.

Una segunda línea es corregir en su raíz la degradación del nodo de plataforma
del fork ante reinicios de servicio repetidos, cuyo síntoma quedó cuantificado en
este trabajo pero cuya causa interna no se investigó. Esta corrección
beneficiaría a cualquier uso intensivo del servicio de reinicio, no solo al
aprendizaje por refuerzo.

En el plano del aprendizaje, los resultados sugieren dos refinamientos: revisar
la bonificación terminal para que la orientación final deficiente no resulte
rentable (el error de yaw residual de 0,44 rad frente a los 0,01 rad del PID
tiene origen en el diseño de la recompensa), y estudiar el escalado del
curriculum hacia trayectorias con múltiples puntos de paso y referencias
dinámicas. Sobre esa base, la extensión jerárquica planteada en los objetivos
(un segundo controlador que transforme velocidades deseadas en tasas angulares)
seguiría siendo la continuación natural hacia el control de bajo nivel.

Finalmente, la comparación con el controlador PID puede ampliarse en dos
direcciones que delimitarían mejor el espacio donde el aprendizaje aporta
ventaja: escenarios con perturbaciones (viento, ruido de estimación, retardos) y
tareas sin ley de control evidente, donde el enfoque clásico requiere diseño
manual específico y la política puede aprenderse con la misma infraestructura sin
cambios.
